\section{如果没有一级标题，那么这就是一级标题}

\subsection{这是二级标题}

\subsubsection{这是三级标题}

\paragraph{六级标题将会变成paragraph}

这是一段正常的文字，\textbf{这里是粗体}，\textit{这里是斜体}。

可以使用Latex的相关行内公式，如$a^2+b^2=c^2$。

可以使用行内代码，但是需要注意的是行内代码不区分代码语言，如 \mintinline{text}|print("Hello World")|。

可以使用超链接功能，如\href{https://www.cleversmall.com}{Cleversmall}。

可以使用文献引用功能，但必须配合\textbf{bibtex}使用，该功能在markdown中只起到一个标志的作用，转为latex后才会生效，如文献\cite{name2024paper}

可以使用图片插入和图片引用跳转功能，如图\ref{fig1}所示。

\begin{figure}[ht]
    \centering
    \includegraphics[width=0.50\textwidth]{./figure/latex_bird.png}
    \label{fig1}
\end{figure}

可以插入表格，并识别对齐方式，并使用引用的跳转，如表\ref{tab1}所示。

\begin{table}[ht]
    \centering
    \caption{测试表格}
    \begin{tabular}{rcl}
        \toprule
        \textbf{标题1}&\textbf{这里居中对齐}&\textbf{标题2} \\
        \midrule
          右对齐了  &     内容2     &  左对齐了   \\
                    &     内容3     &             \\
        \bottomrule
    \end{tabular}
    \label{tab1}
\end{table}

可以使用无序列表功能，如下。

\begin{itemize}
    \item 一个项目
    \item 两个项目
    \item 三个项目
\end{itemize}

也可以使用有序列表功能，如下。

\begin{enumerate}
    \item 内容一
    \item 内容二
    \item 内容三
\end{enumerate}

可以使用行间公式功能，但请使用标准latex公式格式书写其中内容。

\begin{equation}
    \begin{bmatrix}
    a & b \\ c & d
    \end{bmatrix}
    \neq
    \begin{matrix}
    e & f \\ g & h
    \end{matrix}
    \label{eq:1}
\end{equation}

\begin{equation}
    a<b, c>d, e\in (0,1)
\end{equation}

可以使用行间代码功能，并识别语言。

\begin{minted}{python}
import numpy

print("Hello World")
\end{minted}

还有好多功能理论上，后面慢慢加吧

可以使用新的内联 HTML 支持：如下 \underline{下划线文本} 与 \textcolor{Red}{红色文本} 示例。其中数学 $x<y$ 不会被误判。

